\section{Methodology}

% Plan:
% - dataset section: the dataset we used (UCI HAR), mention how we use RAW signals for our analysis
% - corruption framework: introduce our corruption framework with the corruption type, granularity,
%   and severity, subsection each of these (expand on the different types, etc)
% - models: models that we used for testing, why did we pick these, etc

\subsection{Dataset}
Our investigation was focused on the Human Activity Recognition using Smartphones dataset from
Reyes-Ortiz et al. \cite{reyes-ortiz-2013}. The dataset provided was already split between training
and test data. The dataset contains the raw signals as well as a dataset containing 561 features
that were extracted from each signal window using the frequency domain. Due to the feature dataset
removing the time component of the signal, we decided against using the feature dataset because
manipulating the featured signals would not be possible. Our study focuses on the raw signal data,
therefore we applied our corruption framework directly on the raw signal data. Note that it would be
possible to reconstruct the frequency-based feature dataset after a corruption is applied by
performing the feature engineering outlined in their work. However that is beyond the scope
of this study, we simply decided to work against the raw signal data itself and compare against
clean data.

\subsection{Corruption Framework}
The corruption framework is built along three dimensions: \textit{Type}, \textit{Granularity},
and \textit{Severity}.

\subsubsection{Corruption Type}
We investigated various different corruption types:

\begin{enumerate}
    \item Stochastic: Adds random noise to the signal window, simulating stochastic sensor noise.
    \item Dropout: Randomly zeros out a fraction of timesteps in the signal window. Simulates
        intermittent sensor failure or signal loss.
    \item Bias: Adds a constant offset to the signal window. Simulates miscalibration or aging.
    \item Amplification (gain > 1): Scales the signal by a constant multiplier greater than 1.
        Simulates sensor sensitivity amplification.
    \item Attenuation (0 < gain < 1): Scales the signal by a constant multiplier smaller than 1.
        Simulates sensor sensitivity degradation.
    \item Drift: Adds a linearly growing offset over the signal window. Simulates gradual sensor
        drift over time.
\end{enumerate}

\subsubsection{Granularity}
Corruption can be applied at different levels of granularity by specifying which channels are
affected by the corruption. The dataset contains 6 channels across 2 sensors. Body accelerometer
(x, y, z) and gyroscope (x, y, z). This allows corruption to tb eapplied to an entire sensor
modality, a single axis within a modality, or any combination, enabling targeted investigation.

Note that although possible, we decided against investigating applying multiple corruption types
to the data at a time. The reason for this being that we wanted to targeted investigation of
corruption types, mixing corruption types would make harder to attribute a corruption type.

\subsubsection{Severity}
Each corruption type has an associated severity parameter that controls the degree of corruption
applied. The meaning of severity is type-specific.

\begin{enumerate}
    \item Stochastic: Corresponds to the fraction of channel's standard deviation used as the noise
        scale.
    \item Dropout: Corresponds to the fraction of timesteps zeroed out in the window.
    \item Bias: Corresponds to a fraction of the channel's standard deviation added as a constant
        offset.
    \item Amplification: Corresponds to a constant multiplier greater than 1 applied to the signal.
    \item Attenuation: Corresponds to a constant multiplier smaller than 1 applied to the signal.
    \item Drift: Corresponds to a multiplier of the channel's standard deviation by the end of the
        signal window.
\end{enumerate}


\subsection{Models}

We evaluate three classical machine learning models: K-Nearest Neighbors (KNN), Logistic
Regression, and Support Vector Machine (SVM) as well as a more complex model: Long Short-Term
Memory (LSTM) model.

\paragraph{K-Nearest Neighbors (KNN).}
We use a KNN classifier with $k=5$ and Euclidean distance. The model operates directly on flattened
feature vectors without additional preprocessing.

\paragraph{Logistic Regression.}
We implement Logistic Regression using a pipeline that includes feature standardization followed by
a linear probabilistic classifier:

\[
\text{Pipeline} = \text{StandardScaler} \rightarrow \text{LogisticRegression}
\]

We use the \texttt{lbfgs} solver with a maximum of 3000 iterations to ensure convergence.

\paragraph{Support Vector Machine (SVM).}
We use a linear SVM (LinearSVC) combined with feature standardization:

\[
\text{Pipeline} = \text{StandardScaler} \rightarrow \text{LinearSVC}
\]

The maximum number of iterations is set to 5000 to ensure stable training.

% \subsection{Training and Evaluation Protocol}

% All models are trained on clean training data and evaluated on both clean and corrupted test data.

% We report:
% \begin{itemize}
%     \item Accuracy
%     \item Macro F1-score
%     \item Confusion Matrix (for selected experiments)
% \end{itemize}

% For corruption experiments, models are trained once on clean data and reused across all
% corruption severities to isolate the effect of test-time corruption.
